% Preambulo compartido para clases en formato beamer
\documentclass[9pt, xcolor=svgnames]{beamer}
\usepackage[utf8]{inputenc}
\usepackage[T1]{fontenc}
\usepackage[spanish,es-nodecimaldot]{babel}

\usetheme[progressbar=frametitle]{metropolis}
\metroset{sectionpage=none, subsectionpage=none}
\usepackage{appendixnumberbeamer}
\usepackage{booktabs}
\usepackage{amsmath,amssymb,bm,mathtools}
\usepackage{physics}
\usepackage{graphicx}
\usepackage{multicol}
\usepackage{tikz}
\usepackage{minted}
\usemintedstyle{tango}
\setminted{fontsize=\scriptsize, breaklines=true, linenos=true, bgcolor=LightGray!20, frame=lines}
\newminted{python}{fontsize=\scriptsize, breaklines=true, linenos=true, bgcolor=LightGray!20, frame=lines}
\usepackage{pgfplots}
\pgfplotsset{compat=1.18}
\usepackage{ragged2e}
\usepackage{xspace}
\newcommand{\themename}{\textbf{\textsc{metropolis}}\xspace}

\setbeamercolor{block title}{bg=MidnightBlue!85, fg=white}
\setbeamercolor{block body}{bg=MidnightBlue!5, fg=black}
\setbeamercolor{alerted text}{fg=Crimson}
\setbeamercolor{frametitle}{bg=MidnightBlue!10, fg=black}
\setbeamertemplate{blocks}[rounded][shadow=true]


\weektitle{10}{redes densas para senales fisicas y astronomicas}
\titlegraphic{}
\title[PCFI261 -- Semana 10]{\Large Modelos Computacionales PCFI261\\Semana 10: redes densas para senales\\fisicas y astronomicas}
\newcommand{\compactcode}{\fontsize{5}{4.7}\selectfont}

\begin{document}

\begin{frame}[plain]
\centering
\vfill
{\Large\bfseries Modelos Computacionales PCFI261}\\[0.25cm]
{\large Semana 10: redes densas para senales\\fisicas y astronomicas}\\[0.7cm]
{\large Joaquin Peralta}\\[0.15cm]
{\small \texttt{joaquin.peralta@unab.cl}}\\[0.35cm]
{\small Universidad Andres Bello}
\vfill
\end{frame}

\begin{frame}{Organizacion de la semana}
\begin{block}{Continuidad desde Semana 09}
Venimos de redes densas para datos tabulados y del caso Pima Indians Diabetes. Esta semana mantenemos redes densas, escalamiento, train/validation/test y metricas, pero cambiamos el contexto hacia problemas cientificos simulados.
\end{block}

\begin{columns}[T,onlytextwidth]
\column{0.48\textwidth}
\textbf{Clase 1}
\begin{itemize}
    \item emulador para un oscilador amortiguado;
    \item regresion con una red densa;
    \item error fisico e interpolacion;
    \item primer benchmark CPU/GPU.
\end{itemize}

\column{0.48\textwidth}
\textbf{Clase 2}
\begin{itemize}
    \item curvas de luz sinteticas;
    \item transitos de exoplanetas;
    \item clasificacion binaria con red densa;
    \item aumento de escala para justificar GPU.
\end{itemize}
\end{columns}
\end{frame}

\begin{frame}{Idea de la semana}
\begin{columns}[T,onlytextwidth]
\column{0.50\textwidth}
\begin{block}{No queremos una red mas compleja}
Queremos una pregunta cientifica mas interesante usando la misma herramienta: capas densas, activaciones, perdida, optimizador y validacion.
\end{block}
\begin{itemize}
    \item Clase 1: aprender una funcion fisica continua.
    \item Clase 2: reconocer una firma astronomica simple.
    \item En ambas: comparar CPU/GPU sin convertir la clase en ingenieria de hardware.
\end{itemize}
\column{0.46\textwidth}
\centering
\begin{tikzpicture}[scale=0.74,transform shape]
\draw[fill=MidnightBlue!7,draw=MidnightBlue!70,very thick,rounded corners] (-3.0,-2.1) rectangle (3.0,2.1);
\node[font=\small\bfseries,MidnightBlue] at (0,1.75) {Semana 10: misma red, preguntas cientificas};
\draw[very thick,ForestGreen] plot[smooth] coordinates {(-2.55,0.6)(-2.0,1.05)(-1.45,0.35)(-0.85,-0.75)(-0.25,-0.25)(0.35,0.75)(0.95,0.25)(1.55,-0.65)(2.2,-0.2)};
\draw[very thick,Crimson] (-2.45,-1.35) -- (-1.15,-1.35) -- (-1.0,-1.62) -- (-0.35,-1.62) -- (-0.2,-1.35) -- (2.35,-1.35);
\node[circle,fill=orange!80,draw=white,minimum size=13pt] at (-1.3,-0.55) {};
\node[circle,fill=MidnightBlue!85,draw=white,minimum size=23pt] at (-1.0,-0.45) {};
\foreach \x/\y in {-2.1/1.02,-1.6/0.91,-1.1/1.25,-0.6/1.22,-0.1/1.05,0.4/0.92,0.9/1.18,1.4/1.24,1.9/1.03} {\node[circle,fill=Purple!55,draw=white,inner sep=1.4pt] at (\x,\y) {};}
\node[font=\scriptsize,align=center] at (-1.2,-1.95) {oscilador\\regresion};
\node[font=\scriptsize,align=center] at (1.45,-1.95) {transito\\clasificacion};
\end{tikzpicture}

\end{columns}
\end{frame}

\section{Clase 1}

\begin{frame}[plain]
\centering
\vfill
{\Huge \bfseries Clase 1}\\[0.4cm]
{\Large Redes densas como emuladores de modelos fisicos}
\vfill
\end{frame}

\begin{frame}{Objetivos de la Clase 1}
\begin{itemize}
    \item Construir datos sinteticos de un oscilador amortiguado con parametros variables.
    \item Entrenar una red densa para aproximar $x(t; A,\gamma,\omega)$.
    \item Evaluar el error del emulador en curvas que no fueron usadas para entrenar.
    \item Medir tiempos de entrenamiento en CPU/GPU cuando aumenta el tamano del problema.
\end{itemize}
\end{frame}

\begin{frame}{Aplicacion fisica: oscilador amortiguado}
\begin{columns}[T,onlytextwidth]
\column{0.52\textwidth}
\justifying
Un oscilador amortiguado simple puede escribirse, en una version docente, como
\[
 x(t)=A\,e^{-\gamma t}\cos(\omega t+\phi).
\]
En lugar de resolver la ecuacion cada vez, entrenamos una red que aprende un \textbf{emulador}: recibe parametros y tiempo, y devuelve una posicion aproximada.

\begin{block}{Lectura fisica}
La red no descubre la ley desde cero: aprende una aproximacion numerica a partir de ejemplos generados por un modelo conocido.
\end{block}

\column{0.44\textwidth}
\centering
\begin{tikzpicture}[x=0.62cm,y=1.05cm]
\draw[->,thick] (0,-1.25) -- (8.2,-1.25) node[right,font=\scriptsize] {$t$};
\draw[->,thick] (0,-1.25) -- (0,1.45) node[above,font=\scriptsize] {$x(t)$};
\draw[very thick,MidnightBlue] plot[smooth,samples=100,domain=0:7.7] (\x,{exp(-0.18*\x)*cos(150*\x)});
\draw[dashed,ForestGreen,thick] plot[smooth,samples=80,domain=0:7.7] (\x,{exp(-0.18*\x)});
\draw[dashed,ForestGreen,thick] plot[smooth,samples=80,domain=0:7.7] (\x,{-exp(-0.18*\x)});
\node[font=\scriptsize,MidnightBlue] at (4.8,0.85) {senal fisica};
\node[font=\scriptsize,ForestGreen] at (5.7,-0.75) {envolvente};
\end{tikzpicture}

\end{columns}
\end{frame}

\begin{frame}{De la ecuacion a una tabla de entrenamiento}
\centering
\begin{tikzpicture}[>=latex,scale=0.86,transform shape]
\tikzstyle{box}=[draw,rounded corners,very thick,minimum width=2.8cm,minimum height=1.0cm,align=center,font=\small\bfseries]
\node[box,fill=MidnightBlue!12] (params) at (0,0) {parametros\\$A,\gamma,\omega,\phi$};
\node[box,fill=ForestGreen!12] (time) at (3.4,0) {tiempos\\$t$};
\node[box,fill=orange!14] (model) at (6.8,0) {modelo fisico\\$x(t)$};
\node[box,fill=Purple!12] (table) at (10.2,0) {dataset\\$X,y$};
\draw[->,very thick] (params) -- (time);
\draw[->,very thick] (time) -- (model);
\draw[->,very thick] (model) -- (table);
\node[font=\small] at (5.1,-1.25) {Cada fila contiene tiempo + parametros como entrada, y posicion como salida.};
\end{tikzpicture}
\end{frame}

\begin{frame}[fragile]{Parte 1: generar datos fisicos}
\inputminted[fontsize=\compactcode]{python}{codes/clase1-p1.py}
\end{frame}

\begin{frame}{Red densa para regresion}
\begin{columns}[T,onlytextwidth]
\column{0.52\textwidth}
\begin{itemize}
    \item Entrada: $t$, $A$, $\gamma$, $\omega$ y $\phi$.
    \item Salida: posicion $x(t)$.
    \item Perdida: error cuadratico medio, porque la salida es continua.
    \item Activacion final: lineal, no sigmoide.
\end{itemize}
\begin{alertblock}{Cambio respecto a Pima}
Ya no buscamos una probabilidad binaria. Ahora buscamos aproximar una cantidad fisica real.
\end{alertblock}
\column{0.44\textwidth}
\centering
\begin{tikzpicture}[scale=0.66,transform shape]
\def\xA{0} \def\xB{2.6} \def\xC{5.2} \def\xD{7.7}
\foreach \y/\lab in {-1.2/$t$,-0.6/$A$,0/$\gamma$,0.6/$\omega$,1.2/$\phi$} {\node[circle,fill=MidnightBlue!75,draw=white,minimum size=8pt] at (\xA,\y) {};\node[font=\scriptsize] at (\xA-0.45,\y) {\lab};}
\foreach \y in {-1.5,-0.9,-0.3,0.3,0.9,1.5} {\node[circle,fill=ForestGreen!70,draw=white,minimum size=8pt] at (\xB,\y) {};}
\foreach \y in {-1.2,-0.6,0,0.6,1.2} {\node[circle,fill=ForestGreen!70,draw=white,minimum size=8pt] at (\xC,\y) {};}
\node[circle,fill=Crimson!82,draw=white,minimum size=9pt] at (\xD,0) {};
\foreach \ya in {-1.2,-0.6,0,0.6,1.2} {\foreach \yb in {-1.5,-0.9,-0.3,0.3,0.9,1.5} {\draw[gray!35,thin] (\xA+0.12,\ya)--(\xB-0.12,\yb);}}
\foreach \ya in {-1.5,-0.9,-0.3,0.3,0.9,1.5} {\foreach \yb in {-1.2,-0.6,0,0.6,1.2} {\draw[gray!35,thin] (\xB+0.12,\ya)--(\xC-0.12,\yb);}}
\foreach \ya in {-1.2,-0.6,0,0.6,1.2} {\draw[gray!45,thin] (\xC+0.12,\ya)--(\xD-0.12,0);}
\node[font=\small\bfseries] at (\xA,-2.1) {entrada};
\node[font=\small\bfseries] at (\xB,-2.1) {densa};
\node[font=\small\bfseries] at (\xC,-2.1) {densa};
\node[font=\small\bfseries] at (\xD,-2.1) {$\hat{x}(t)$};
\end{tikzpicture}

\end{columns}
\end{frame}

\begin{frame}[fragile]{Parte 2: entrenar el emulador}
\inputminted[fontsize=\compactcode]{python}{codes/clase1-p2.py}
\end{frame}

\begin{frame}{Que significa que el emulador funcione}
\begin{itemize}
    \item Debe reproducir amplitud, fase y decaimiento en curvas no vistas.
    \item El error no debe concentrarse solo en una zona temporal.
    \item Si extrapolamos fuera del rango de parametros entrenado, el modelo puede fallar con confianza aparente.
\end{itemize}

\begin{block}{Pregunta para el curso}
Si conocemos la ecuacion, por que usar una red? Porque en problemas reales el simulador puede ser caro, ruidoso o no tener inversa simple.
\end{block}
\end{frame}

\begin{frame}[fragile]{Parte 3: evaluar una curva nueva}
\inputminted[fontsize=\compactcode]{python}{codes/clase1-p3.py}
\end{frame}

\begin{frame}{CPU versus GPU: cuando empieza a importar}
\begin{columns}[T,onlytextwidth]
\column{0.52\textwidth}
\begin{itemize}
    \item Para miles de ejemplos, una CPU puede bastar.
    \item La GPU ayuda cuando hay muchas operaciones matriciales similares.
    \item Aumentar el dataset sintetico permite medir el cambio sin buscar archivos externos.
    \item La GPU acelera entrenamiento, pero no arregla mala validacion ni mala pregunta fisica.
\end{itemize}
\column{0.44\textwidth}
\centering
\begin{tikzpicture}[scale=0.78,transform shape]
\draw[->,thick] (0,0) -- (5.6,0) node[right,font=\scriptsize] {tamano};
\draw[->,thick] (0,0) -- (0,3.5) node[above,font=\scriptsize] {tiempo};
\draw[very thick,Crimson] plot[smooth] coordinates {(0.3,0.35)(1.3,0.9)(2.4,1.55)(3.5,2.35)(4.8,3.15)};
\draw[very thick,MidnightBlue] plot[smooth] coordinates {(0.3,0.55)(1.3,0.72)(2.4,0.95)(3.5,1.18)(4.8,1.45)};
\draw[dashed,gray,thick] (1.8,0) -- (1.8,3.1);
\node[Crimson,font=\scriptsize\bfseries] at (4.15,3.15) {CPU};
\node[MidnightBlue,font=\scriptsize\bfseries] at (4.25,1.05) {GPU};
\node[font=\scriptsize,align=center] at (1.25,3.15) {problema\\pequeno};
\node[font=\scriptsize,align=center] at (3.25,0.35) {paralelismo\\visible};
\end{tikzpicture}

\end{columns}
\end{frame}

\begin{frame}[fragile]{Parte 4: benchmark simple CPU/GPU}
\inputminted[fontsize=\compactcode]{python}{codes/clase1-p4.py}
\end{frame}

\begin{frame}{Cierre Clase 1}
\begin{itemize}
    \item Una red densa puede actuar como emulador de una funcion fisica parametrizada.
    \item En regresion cambia la salida, la perdida y la interpretacion del error.
    \item La comparacion CPU/GPU debe hacerse con problemas suficientemente grandes y con metricas claras.
\end{itemize}
\end{frame}

\section{Clase 2}

\begin{frame}[plain]
\centering
\vfill
{\Huge \bfseries Clase 2}\\[0.4cm]
{\Large Redes densas para detectar transitos de exoplanetas}
\vfill
\end{frame}

\begin{frame}{Objetivos de la Clase 2}
\begin{itemize}
    \item Simular curvas de luz simples con y sin transito planetario.
    \item Entrenar una red densa sobre vectores de flujo, sin usar CNN.
    \item Evaluar matriz de confusion, precision, recall y umbral de decision.
    \item Escalar el experimento para observar diferencias de rendimiento con GPU.
\end{itemize}
\end{frame}

\begin{frame}{Aplicacion astronomica: una caida de brillo}
\begin{columns}[T,onlytextwidth]
\column{0.52\textwidth}
\justifying
Cuando un planeta pasa frente a su estrella, el brillo observado puede bajar levemente durante un intervalo corto. En esta clase generamos curvas sinteticas con ruido y entrenamos una red densa para distinguir:
\begin{itemize}
    \item estrella sin transito;
    \item estrella con una senal de transito simple.
\end{itemize}

\begin{alertblock}{Limite del modelo}
No haremos astronomia de precision. Buscamos una experiencia docente: senal, ruido, clasificacion y validacion.
\end{alertblock}

\column{0.44\textwidth}
\centering
\begin{tikzpicture}[scale=0.72,transform shape]
\shade[inner color=orange!70,outer color=orange!20] (-1.55,0.75) circle (1.05);
\node[circle,fill=MidnightBlue!85,draw=white,minimum size=20pt] at (-1.15,0.92) {};
\draw[thick,->] (-3.0,0.75) -- (0.6,0.75);
\draw[->,thick] (-3.0,-1.55) -- (2.8,-1.55) node[right,font=\scriptsize] {$t$};
\draw[->,thick] (-3.0,-1.55) -- (-3.0,-0.15) node[above,font=\scriptsize] {flujo};
\draw[very thick,MidnightBlue] (-2.85,-0.45) -- (-1.15,-0.45) -- (-0.9,-0.92) -- (-0.15,-0.92) -- (0.1,-0.45) -- (2.45,-0.45);
\node[font=\scriptsize] at (0.3,1.75) {transito planetario};
\node[font=\scriptsize,MidnightBlue] at (1.35,-0.85) {caida de brillo};
\end{tikzpicture}

\end{columns}
\end{frame}

\begin{frame}{Por que todavia no usamos CNN}
\begin{itemize}
    \item Una curva de luz discretizada puede verse como un vector: $[f_1,f_2,\ldots,f_N]$.
    \item Una red densa puede aprender combinaciones globales de esos puntos.
    \item Las CNN seran utiles mas adelante para explotar estructura local, ventanas y patrones desplazados.
    \item Aqui mantenemos el foco en redes densas, validacion y GPU.
\end{itemize}
\end{frame}

\begin{frame}[fragile]{Parte 1: simular curvas de luz}
\inputminted[fontsize=\compactcode]{python}{codes/clase2-p1.py}
\end{frame}

\begin{frame}{Arquitectura del clasificador}
\centering
\begin{tikzpicture}[scale=0.67,transform shape]
\def\xA{0} \def\xB{2.8} \def\xC{5.6} \def\xD{8.3}
\foreach \y in {-1.8,-1.2,-0.6,0,0.6,1.2,1.8} {\node[circle,fill=MidnightBlue!75,draw=white,minimum size=7pt] at (\xA,\y) {};}
\foreach \y in {-2.0,-1.4,-0.8,-0.2,0.4,1.0,1.6,2.2} {\node[circle,fill=ForestGreen!70,draw=white,minimum size=7pt] at (\xB,\y) {};}
\foreach \y in {-1.5,-0.9,-0.3,0.3,0.9,1.5} {\node[circle,fill=ForestGreen!70,draw=white,minimum size=7pt] at (\xC,\y) {};}
\node[circle,fill=Crimson!82,draw=white,minimum size=9pt] at (\xD,0) {};
\foreach \ya in {-1.8,-1.2,-0.6,0,0.6,1.2,1.8} {\foreach \yb in {-2.0,-1.4,-0.8,-0.2,0.4,1.0,1.6,2.2} {\draw[gray!30,thin] (\xA+0.12,\ya)--(\xB-0.12,\yb);}}
\foreach \ya in {-2.0,-1.4,-0.8,-0.2,0.4,1.0,1.6,2.2} {\foreach \yb in {-1.5,-0.9,-0.3,0.3,0.9,1.5} {\draw[gray!30,thin] (\xB+0.12,\ya)--(\xC-0.12,\yb);}}
\foreach \ya in {-1.5,-0.9,-0.3,0.3,0.9,1.5} {\draw[gray!45,thin] (\xC+0.12,\ya)--(\xD-0.12,0);}
\node[font=\small\bfseries,align=center] at (\xA,-2.65) {curva de luz\\128 puntos};
\node[font=\small\bfseries] at (\xB,-2.65) {64};
\node[font=\small\bfseries] at (\xC,-2.65) {32};
\node[font=\small\bfseries,align=center] at (\xD,-2.65) {$p$\\transito};
\end{tikzpicture}

\end{frame}

\begin{frame}[fragile]{Parte 2: entrenar una red densa}
\inputminted[fontsize=\compactcode]{python}{codes/clase2-p2.py}
\end{frame}

\begin{frame}{Metricas: mirar errores, no solo accuracy}
\begin{columns}[T,onlytextwidth]
\column{0.50\textwidth}
\centering
\begin{tikzpicture}[scale=0.82,transform shape]
\foreach \x/\y/\v/\lab/\shade in {0/1.6/240/TN/45,2.1/1.6/18/FP/18,0/0/26/FN/22,2.1/0/216/TP/36} {
  \draw[fill=MidnightBlue!\shade,draw=white,very thick] (\x,\y) rectangle +(1.9,1.35);
  \node[font=\Large\bfseries] at (\x+0.95,\y+0.82) {\v};
  \node[font=\small] at (\x+0.95,\y+0.35) {\lab};
}
\node[font=\small] at (0.95,-0.35) {pred no};
\node[font=\small] at (3.05,-0.35) {pred si};
\node[font=\small,rotate=90] at (-0.45,2.25) {real no};
\node[font=\small,rotate=90] at (-0.45,0.65) {real si};
\end{tikzpicture}
\column{0.46\textwidth}
\begin{itemize}
    \item Falsos positivos: candidatos que no son transitos.
    \item Falsos negativos: transitos que el modelo pierde.
    \item El umbral decide que tan conservador o sensible queremos ser.
\end{itemize}
\end{columns}
\end{frame}

\begin{frame}[fragile]{Parte 3: evaluar el clasificador}
\inputminted[fontsize=\compactcode]{python}{codes/clase2-p3.py}
\end{frame}

\begin{frame}{Actividad breve: cambiar la dificultad}
\begin{block}{Experimentos sugeridos}
Cambiar un parametro a la vez y observar el efecto en validacion:
\end{block}
\begin{itemize}
    \item aumentar el ruido de las curvas;
    \item reducir la profundidad del transito;
    \item mover aleatoriamente el centro del transito;
    \item cambiar el numero de puntos temporales;
    \item comparar una red pequena contra una red mas ancha.
\end{itemize}
\end{frame}

\begin{frame}[fragile]{Parte 4: experimento de escala para GPU}
\inputminted[fontsize=\compactcode]{python}{codes/clase2-p4.py}
\end{frame}

\begin{frame}{Discusion final}
\begin{itemize}
    \item Que diferencia hay entre aprender una funcion fisica y clasificar una senal astronomica?
    \item Como cambia la interpretacion entre MSE, accuracy, recall y falsos positivos?
    \item En que punto la GPU se vuelve pedagogicamente visible?
    \item Que informacion local aprovecharia una CNN cuando veamos imagenes o series mas adelante?
\end{itemize}
\end{frame}

\begin{frame}{Cierre de la semana}
\begin{itemize}
    \item Las redes densas sirven para mas que salud o datos tabulados: tambien pueden aproximar modelos y reconocer senales cientificas simples.
    \item La complejidad debe crecer desde la pregunta, no desde la arquitectura.
    \item La GPU es una herramienta de escala: acelera calculos matriciales grandes, pero no reemplaza validacion ni interpretacion fisica.
\end{itemize}
\end{frame}

\end{document}
